\subsection{The relative lifting argument and the actual CM witness}

\begin{result}{E.46}{Proved}{Relative peripheral lifting and the full gentle witness}
Let $\Gamma$ be $\widehat F(x_1,x_2)$ or
$C_2(s)\amalg C_3(r)$, put $x_3=(x_1x_2)^{-1}$,
and use $x_1=s,x_2=r$ in the second case.
For an open normal subgroup $U$, let $R$ be its closed
nilpotent residual and put $\Gamma_{\rm rel}=\Gamma/R$.
Every continuous automorphism $\alpha$ of this actual
quotient with specified conjugators
$\alpha(\bar x_i)=b_i\bar x_i b_i^{-1}$ lifts to an
automorphism $\Phi$ of $\Gamma$ with conjugators lifting
all three $b_i$ and
\[
 q_{\rm rel}\Phi=\alpha q_{\rm rel}.
                                                               \tag{E.232}
\]
If one $b_i=1$, the lift can fix that peripheral literally.
Here peripheral preservation has exponent one, as in the
displayed equation.

In particular the element designated in LM7 E.58 exists
in the full group $\widehat{\mathrm{GT}}_0$ with its
specified genuine frames. It has multiplier one, trivial
peripheral homology, and the CM parity $(1,0)$ of E.24.
It therefore lies outside the arithmetic image.
E.49 proves that the same CM character vanishes on
all of GT, and consequently also excludes this
witness from GT. The equality of the arithmetic
image with GT remains the separate question in E.24.
\end{result}
\begin{proof}
We supply the finite counting argument used by the
relative compactness step. This independently replaces
the unavailable original source D.50.5--17; the referee's
verdict is not used as a mathematical input.

\medskip\noindent\textbf{The precise finite assertion.}
Let $\Lambda$ be the corresponding abstract source and
$q_0:\Lambda\twoheadrightarrow E$ a finite quotient.
Fix $A\triangleleft E$, write $D=E/A$, let
$\chi_0:\Lambda\to D$ be the quotient, and put
$X_i=q_0(x_i)$ and $T_i=X_i^A$.
We construct a finite central extension
$p:D^\dagger\to D$ and a surjective actual source map
$q_0^\dagger:\Lambda\to D^\dagger$ over $\chi_0$.
Given another epimorphism $q_1^\dagger$ and elements
$h_i^\dagger$ satisfying
$q_1^\dagger(x_i)=h_i^\dagger q_0^\dagger(x_i)
(h_i^\dagger)^{-1}$, choose $H_i\in E$ above
$p(h_i^\dagger)$ and put $T_i^1=H_iT_iH_i^{-1}$,
$T_i^0=T_i$.
We prove equality of the numbers of epimorphisms
$q_j:\Lambda\to E$ with $q_j(x_i)\in T_i^j$,
for $j=0,1$. The reference map is counted on the
first side, so the second side is positive.
Membership in $T_i^1$ then gives a conjugator
$H_i a_i$, $a_i\in A$, with its prescribed
$D$-coordinate. No hypothesis $[E,A]=A$ is required.

\medskip\noindent\textbf{The finite implementer extensions.}
For every subgroup $J\leq E$ onto $D$ such that all
$T_i\cap J$ are nonempty, put $M=J\cap A$ and
$e=\exp(J)$. For every irreducible complex representation
$\rho:M\to\operatorname{GL}(V)$ whose isomorphism class
is invariant under $J$, put $d=\dim V$ and define
\[
 \widetilde J_\rho=
 \left\{(g,T):\begin{array}{l}
  g\in J,\quad T\rho(m)T^{-1}=\rho(gmg^{-1})\ (m\in M),\\
  (\det T)^e=1
 \end{array}\right\}.
                                                               \tag{E.233}
\]
It is a group. For each $g$ an implementer exists by
invariance, and Schur's lemma says that all implementers
are scalar multiples. The determinant condition leaves
exactly $de$ choices. Thus this is a finite extension
of $J$ by the central scalar group $\mu_{de}$.
The graph $m\mapsto(m,\rho(m))$ is an embedded normal
subgroup: $\rho(m)^e=1$ proves the determinant condition,
and the implementer equation proves normality.
It intersects the scalar kernel trivially.

Quotient by this graph and call the result $Q_\rho$.
There is an exact central extension
$1\to\mu_{de}\to Q_\rho\to D\to1$ and an actual
isomorphism
$\widetilde J_\rho\simeq J\times_DQ_\rho$.
Indeed the map $(g,T)\mapsto(g,[(g,T)])$ is injective,
and both groups have $|J|de$ elements. Projection to
the second matrix coordinate gives a representation
$\Psi_\rho:J\times_DQ_\rho\to\operatorname{GL}(V)$
whose restriction to $M\times\{1\}$ is $\rho$.

For the free source choose any lifts of
$\chi_0(x_1),\chi_0(x_2)$ in $Q_\rho$.
For the modular source choose
$a_s\in T_1\cap J$, $a_r\in T_2\cap J$.
Their orders divide two and three. An implementer for
$a_s$ has scalar square, so it can be rescaled to
$T_s^2=1$; its determinant then has order dividing two
and hence $e$. If $a_s=1$, choose $T_s=1$ instead.
Likewise choose $T_r^3=1$.
The images of these pairs in $Q_\rho$ satisfy the
actual source torsion relations. In either case we
have a homomorphism $R_\rho^0:\Lambda\to Q_\rho$;
its third value is the inverse product of the first
two, not an independent choice.

Take $D^\dagger$ to be the joint image of $\Lambda$
under $\chi_0$ and all these $R_\rho^0$.
It maps onto $D$, with kernel in the product of the
scalar groups, hence central. Its source map
$q_0^\dagger$ is onto by definition. This uses actual
finite maps and inserts no extension splitting.

\medskip\noindent\textbf{The exact pants count.}
Fix $J\leq E$ onto $D$, put $M=J\cap A$, and let
$C_i^j=T_i^j\cap J$ and
$B_i^j=\sum_{a\in C_i^j}a\in\mathbb C[J]$.
Each class sum commutes with $\mathbb C[M]$ and is
supported over the one quotient element
$\chi_j(x_i)$, where $\chi_j=pq_j^\dagger$.
Thus $B_1^jB_2^jB_3^j\in\mathbb C[M]$, and its
identity coefficient is exactly the number $H_j(J)$
of homomorphisms $\Lambda\to J$ with these peripheral
conditions. In the modular case the torsion relations
are automatic on the specified conjugacy classes.

For a primitive central idempotent $e_\rho$ of
$\mathbb C[M]$, if the coset supporting $B_i^j$
moves $[\rho]$, then $e_\rho B_i^j=0$.
To see this, combine the equality
$B_i^je_\rho=e_\rho B_i^j$ with
$a e_\rho=e_{a\rho}a$ for every summand.
The two quotient generators generate $D$, so only
the irreducible classes invariant under all of $J$
can contribute. These are precisely the classes
used above, for either value of $j$.

For such a $\rho$, let $R_\rho^j$ be the
$Q_\rho$-coordinate of $q_j^\dagger$, and put
$L_i^j=\sum_{a\in C_i^j}\Psi_\rho(a,R_\rho^j(x_i))$.
Conjugation by $\rho(m)$ permutes the summands and
leaves their $Q_\rho$-coordinates unchanged.
Thus $L_i^j$ is scalar by irreducibility.
The product of the three $Q_\rho$-coordinates is one,
so
$\rho(B_1^jB_2^jB_3^j)=L_1^jL_2^jL_3^j$.
The regular character of $M$ now gives the exact
formula
\[
 H_j(J)=\frac1{|M|}
   \sum_{\substack{\rho\in\operatorname{Irr}(M)\\
                   [\rho]\text{ invariant under }J}}
       \frac{\operatorname{tr}(L_1^j)
             \operatorname{tr}(L_2^j)
             \operatorname{tr}(L_3^j)}{\dim\rho}.
                                                               \tag{E.234}
\]
In particular the denominator is $\dim\rho$:
each scalar $L_i$ has value $\operatorname{tr}(L_i)/d$,
and the regular-character term is
$d\operatorname{tr}(L_1L_2L_3)$.

Choose $j_i\in J$ above $p(h_i^\dagger)$.
Since $j_i^{-1}H_i\in A$, the exact class-set equality is
$C_i^1=j_iC_i^0j_i^{-1}$.
Emptiness therefore agrees on both sides.
For the $Q_\rho$-coordinate $h_{i,\rho}$ of
$h_i^\dagger$, the pair $(j_i,h_{i,\rho})$ belongs
to $J\times_DQ_\rho$, and the specified source
conjugator identity gives
\[
 L_i^1=\Psi_\rho(j_i,h_{i,\rho})
           L_i^0\Psi_\rho(j_i,h_{i,\rho})^{-1}.
\]
All traces in (E.234) are unchanged, proving
$H_0(J)=H_1(J)$. A subgroup not onto $D$ has
both counts zero. Thus equality holds for every
$J\leq E$. If $P_j(K)$ counts maps with exact
image $K$, then $H_j(J)=\sum_{K\leq J}P_j(K)$.
M\"obius inversion on the finite subgroup poset
gives $P_0(E)=P_1(E)>0$, proving the finite assertion
and its exact conjugator-coordinate conclusion.

\medskip\noindent\textbf{Relative compactness and the local product.}
In the compact space $\Gamma^3$, impose the
peripheral product equation, generation of every
finite quotient, and all prescribed quotient
coordinates of the three conjugators. Every finite
collection of these closed conditions is contained
in one actual finite quotient $q_0:\Gamma\to E_0$
which also retains $\Gamma/U$.
Put $U_0=q_0(U)$, $A_0=\gamma_\infty(U_0)$, and
$D_0=E_0/A_0$.
Every retained relative-coordinate quotient kills
$A_0$, since its image of $U$ is nilpotent.

Apply the finite construction to $E_0,A_0$.
The image of $U$ in its central extension
$D^\dagger\to D_0$ is a finite central extension
of the nilpotent group $U_0/A_0$; it is therefore
nilpotent. Its actual source map factors through
$\Gamma_{\rm rel}$. Precompose by $\alpha$.
The new map is onto and has the literal images
of $b_i$ as its three conjugators. The finite
assertion supplies an epimorphism to $E_0$ with
those exact $D_0$-coordinates. Lifting its three
conjugators to $\Gamma$ meets the finite collection.
Compactness now gives one triple meeting every
condition, hence an endomorphism $\Phi$.
Its compact image surjects onto every finite
quotient, so it equals $\Gamma$.

For completeness, such an epimorphism is injective.
Preimage under a surjection is an injection on
the finite set of open subgroups of any fixed
index, hence permutes it. Every open subgroup
is consequently a preimage and contains the
kernel. Their intersection is trivial.
This is the required Hopficity argument.
The retained coordinates give (E.232) exactly.
If $b_i=1$, its lift lies in $R$; conjugating
by that lift's inverse fixes the peripheral
literally and changes none of the retained
coordinates. Notice that $\alpha$ need not
preserve $\Gamma/U$: the finite count allowed
both quotient maps $\chi_0$ and $\chi_1$.

Let $O^p(U)$ be the kernel of the maximal pro-$p$
quotient. Sylow decomposition of finite nilpotent
groups and coprime Goursat give
\[
 U/R\simeq\prod_p U^{(p)},\qquad
 \Gamma/R\simeq
       \prod_{p,\,\Gamma/U}\Gamma/O^p(U).
                                                               \tag{E.235}
\]
Here every finite selection of distinct-prime
quotients is jointly onto because their common
quotients have coprime orders; compactness gives
the infinite product. For the second map, lift
the common $\Gamma/U$-coordinate by one element
of $\Gamma$, then lift the discrepancy using
the first map. These are actual source maps,
with no chosen splitting of an extension.

\medskip\noindent\textbf{The specified finite sources.}
Now put
\[
 \Gamma=C_2(s)\amalg C_3(r),\quad \tau=sr,\quad
 F=\ker(\Gamma\to\mathfrak S_3),\quad
 x=\tau^{-2},\quad y=sxs,\quad H=F'F^3.
                                                               \tag{E.236}
\]
The six epimorphisms to $\mathfrak S_3$ have one
kernel, since its automorphism group permutes
them transitively. Thus $F$ and $H$ are
characteristic. The standard index-six free
subgroup has free basis $x,y$; its quotient
graph, or the direct Reidemeister--Schreier
calculation, has six edges and five vertices.
Consequently $[\Gamma:H]=6\cdot3^2=54$.

Use the actual quotient
$q:\Gamma\to Q=\operatorname{SL}_2(\mathbb Z/12)/\{\pm I\}$
with $s$ sent to
$[\left(\begin{smallmatrix}0&1\\-1&0\end{smallmatrix}\right)]$
and $r$ to
$[\left(\begin{smallmatrix}0&-1\\1&-1\end{smallmatrix}\right)]$.
The signed $F$-image is $V_4\times\operatorname{SL}_2(\mathbb F_3)$,
of order $96$, with the four-element group at
modulus four generated by $U^2,L^2$ of E.24.
The two unipotents generate $\operatorname{SL}_2(\mathbb F_3)$.
Its commutator subgroup is $Q_8$: the matrices
$\left(\begin{smallmatrix}0&1\\2&0\end{smallmatrix}\right)$
and $\left(\begin{smallmatrix}1&1\\1&2\end{smallmatrix}\right)$
square to $-I$, anticommute and generate a
normal group of index three, and the actual
unipotent commutator and its conjugates
generate it. Hence this special linear group
has no nontrivial common quotient with $V_4$;
Goursat gives the asserted direct product.
Taking commutators and cubes therefore gives
\[
 \begin{gathered}
 q(H)=V_4\times Q_8,\qquad |q(H)|=32,\qquad
 z_0=q([x,y]^2)=[5I],\\
 J=\operatorname{im}(\Gamma\to Q\times\Gamma/H)
   =Q\times_{Q/q(H)}(\Gamma/H),\qquad |J|=1728.
 \end{gathered}                                                   \tag{E.237}
\]
The joint-image equality follows by lifting
the prescribed $\Gamma/H$-coordinate and
correcting the discrepancy in $q(H)$.
The literal word $[x,y]^2\in H$ supplies
$(z_0,1)\in C_J(r)$, and $z_0$ is central
of order two. The matrix products are also
checked independently in \texttt{E\_CM\_checks.py}.
These maps factor through
$\Gamma_2=\Gamma/O^2(H)$, and the kernel of
$\Gamma_2\to J$ is a closed subgroup of
$H^{(2)}$, hence pro-two.

\medskip\noindent\textbf{The relative centralizer and the full witness.}
Centralizers of an element of finite order
prime to $p$ lift across a surjection with
pro-$p$ kernel. Here is the needed elementary
form of coprime complement conjugacy.
For a finite $p$-group $K$, induct on $|K|$,
using the characteristic elementary abelian
subgroup of order-$p$ elements of $Z(K)$.
After conjugating in the quotient, the
discrepancy lies in this subgroup $A$ and
is a cyclic cocycle. The norm-zero vectors
of $A$ are the image of $a-1$, since the
order of $a$ is invertible in $\mathbb F_p$.
Conjugation by that vector removes the
discrepancy. This proves complement
conjugacy for the finite cyclic coprime
quotient. Apply it in every invariant
finite quotient of a pro-$p$ kernel;
compactness of the closed conjugator sets
gives the same conclusion there.

Specifically, for a lift $g$ of a target
centralizer element, $\langle a\rangle$
and $\langle gag^{-1}\rangle$ are such
complements in the inverse image of
$\langle\bar a\rangle$. Their identical
generator image gives $u\in K$ with
$ugag^{-1}u^{-1}=a$. Thus $ug$ is the
prescribed centralizer lift.
Apply this to $\Gamma_2\to J$, $p=2$,
$a=r$, to obtain
$c\in C_{\Gamma_2}(r)\cap H^{(2)}$
mapping to $(z_0,1)$.

In (E.235), with $U=H$, use the identity
automorphism and the conjugators
$b_s=b_\tau=1$ and
$b_r=(c\text{ at }2,\,1\text{ at every other prime})$.
Their common finite-base coordinate is one,
and $b_r$ centralizes $r$. The already proved
lifting theorem supplies an automorphism
of $\Gamma$. Normalize by its
$\tau$-conjugator in $R_{\rm nil}(H)$.
We obtain actual full frames with
\[
 \begin{gathered}
 \Phi(\tau)=\tau,\qquad
 \Phi(s)=\beta s\beta^{-1},\qquad
 \Phi(r)=ArA^{-1},\\
 \beta\in R_{\rm nil}(H),\quad A\in H,\quad
 q(\beta)=1,\quad q(A)=[5I],\quad q\Phi=q.
 \end{gathered}                                                   \tag{E.238}
\]
The third peripheral used in the lifting
theorem is $\tau^{-1}$; it has the same
conjugator as $\tau$.

Put $f=s\Phi(s)$. Since the relative action
is the identity,
$f\in R_{\rm nil}(H)\subset H'\subset F'$.
The source relations give
\[
 \begin{gathered}
 \Phi(s)=sf,\quad \Phi(r)=f^{-1}r,\quad
 \Phi(x)=x,\quad \Phi(y)=f^{-1}yf,\\
 f(y,x)f(x,y)=1,\qquad
 f(z,x)f(y,z)f(x,y)=1,
       \quad z=(xy)^{-1}.
 \end{gathered}                                                   \tag{E.239}
\]
Indeed conjugation by $s$ interchanges $x,y$,
and conjugation by $r$ cyclically sends
$x$ to $y$ to $z$ to $x$, directly in
$C_2*C_3$. The two cycles in (E.239) are
therefore exactly $\Phi(s)^2=1$ and
$\Phi(r)^3=1$. The characteristic subgroup
$F$ is preserved and the restriction is
an automorphism. The pair $(1,f)$ is
thus an actual element of
$\widehat{\mathrm{GT}}_0$ in the
Harbater--Schneps definition.

The pure frames in (E.238) are unique.
Discretely the centralizers of $s,r$
in $C_2*C_3$ are their finite factors.
Theorem A of \cite{ET_BZ2024}, applied
to this virtually free group at the
primes two and three, gives the same
completed centralizers. Both factors
meet $F$ trivially. The theorem's
cohomological hypotheses hold here:
a finite-index free subgroup is good
and of finite type, and the finite
extension and discrete-module comparisons
are the same ones verified in E.44.

Finally $\beta\in F'$ makes its
abelian frame coordinate $\rho$ zero.
Thus the normalized frames of E.24
are exactly $T=q(A)^{-1}=5I$ and
$W=q(\beta)^{-1}=I$, with no nonzero
shear. The first is $I$ modulo four
and $-I$ modulo three, so its sign
on $\mathbb F_3(1,-1)$ is one;
the second has parity zero. E.24's
independent field and matrix calculation
excludes this actual element from the
arithmetic image. This proves existence
of the specified full gentle witness
and all its genuine finite frames.
E.49's independent three-division quartic
comparison proves that its CM parity
excludes the five-point pentagon and
hence membership in GT.
\end{proof}
